% GENERATED FILE. Edit canonical lessons, then run npm run generate:books.
% canonical-source-hash: fnv1a64:6382fd24797d4970
% canonical-lessons: TE-C256-sandu, TE-C256-sorangam, TE-C256-sivaru, TE-C256-godamu, TE-C256-kortu

\chapter{Lane, Tunnel, Outskirts, Warehouse, Court}
\label{ch:te-lane-tunnel-outskirts-warehouse-court}
\hlchaptermodality{256}

\begin{chapteropening}
\textbf{By the end of this chapter:} \emph{I can say a lane, a tunnel, the outskirts, a warehouse and a court.}

Complete the last lesson of chapter 256: I can say a lane, a tunnel, the outskirts, a warehouse and a court.
\end{chapteropening}

\section[\texorpdfstring{\te{సందు} (sandu)}{సందు (sandu)}]{\te{సందు} (sandu) — a lane, an alley}

\label{lesson:TE-C256-sandu}



\begin{quote}
Before the new one: say the Telugu for a city, then the Telugu for the village meeting platform.
\end{quote}

\subsection*{You'll want to know: \te{సందు}}
\textbf{\te{సందు}} — \emph{sandu} — \textquotedblleft{}a lane, an alley\textquotedblright{}.

\begin{grammarlens}[title={What you've built}]
One more word for this chapter.
\end{grammarlens}

\subsection*{Your turn}
\begin{itemize}
\raggedright
  \item \emph{Say it:} \emph{sandu}
  \item \emph{Say it:} \emph{sandu}, once more
  \item \emph{Say it:} say \emph{sandu}
  \item \emph{Recall:} say the Telugu for a city, then the Telugu for the village meeting platform, then say \emph{sandu} again
\end{itemize}

\begin{tcolorbox}[breakable,colback=teal!4,colframe=teal!35!black,title={Before you move on}]
What does \te{సందు} mean? (\textquotedblleft{}A lane, an alley\textquotedblright{}.) Say it once more.
\end{tcolorbox}
\section[\texorpdfstring{\te{సొరంగం} (soraṅgaṁ)}{సొరంగం (soraṅgaṁ)}]{\te{సొరంగం} (soraṅgaṁ) — a tunnel}

\label{lesson:TE-C256-sorangam}



\begin{quote}
Before the new one: say the Telugu for the village meeting platform, then the Telugu for a lane.
\end{quote}

\subsection*{You'll want to know: \te{సొరంగం}}
\textbf{\te{సొరంగం}} — \emph{soraṅgaṁ} — \textquotedblleft{}a tunnel\textquotedblright{}.

\begin{grammarlens}[title={What you've built}]
One more word for this chapter.
\end{grammarlens}

\subsection*{Your turn}
\begin{itemize}
\raggedright
  \item \emph{Say it:} \emph{soraṅgaṁ}
  \item \emph{Say it:} \emph{soraṅgaṁ}, once more
  \item \emph{Say it:} say \emph{soraṅgaṁ}
  \item \emph{Recall:} say the Telugu for the village meeting platform, then the Telugu for a lane, then say \emph{soraṅgaṁ} again
\end{itemize}

\begin{tcolorbox}[breakable,colback=teal!4,colframe=teal!35!black,title={Before you move on}]
What does \te{సొరంగం} mean? (\textquotedblleft{}A tunnel\textquotedblright{}.) Say it once more.
\end{tcolorbox}
\section[\texorpdfstring{\te{శివారు} (śivāru)}{శివారు (śivāru)}]{\te{శివారు} (śivāru) — the outskirts}

\label{lesson:TE-C256-sivaru}



\begin{quote}
Before the new one: say the Telugu for a lane, then the Telugu for a tunnel.
\end{quote}

\subsection*{You'll want to know: \te{శివారు}}
\textbf{\te{శివారు}} — \emph{śivāru} — \textquotedblleft{}the outskirts\textquotedblright{}.

\begin{grammarlens}[title={What you've built}]
One more word for this chapter.
\end{grammarlens}

\subsection*{Your turn}
\begin{itemize}
\raggedright
  \item \emph{Say it:} \emph{śivāru}
  \item \emph{Say it:} \emph{śivāru}, once more
  \item \emph{Say it:} say \emph{śivāru}
  \item \emph{Recall:} say the Telugu for a lane, then the Telugu for a tunnel, then say \emph{śivāru} again
\end{itemize}

\begin{tcolorbox}[breakable,colback=teal!4,colframe=teal!35!black,title={Before you move on}]
What does \te{శివారు} mean? (\textquotedblleft{}The outskirts\textquotedblright{}.) Say it once more.
\end{tcolorbox}
\section[\texorpdfstring{\te{గోదాము} (gōdāmu)}{గోదాము (gōdāmu)}]{\te{గోదాము} (gōdāmu) — a warehouse}

\label{lesson:TE-C256-godamu}



\begin{quote}
Before the new one: say the Telugu for a tunnel, then the Telugu for the outskirts.
\end{quote}

\subsection*{You'll want to know: \te{గోదాము}}
\textbf{\te{గోదాము}} — \emph{gōdāmu} — \textquotedblleft{}a warehouse\textquotedblright{}.

\begin{grammarlens}[title={What you've built}]
One more word for this chapter.
\end{grammarlens}

\subsection*{Your turn}
\begin{itemize}
\raggedright
  \item \emph{Say it:} \emph{gōdāmu}
  \item \emph{Say it:} \emph{gōdāmu}, once more
  \item \emph{Say it:} say \emph{gōdāmu}
  \item \emph{Recall:} say the Telugu for a tunnel, then the Telugu for the outskirts, then say \emph{gōdāmu} again
\end{itemize}

\begin{tcolorbox}[breakable,colback=teal!4,colframe=teal!35!black,title={Before you move on}]
What does \te{గోదాము} mean? (\textquotedblleft{}A warehouse\textquotedblright{}.) Say it once more.
\end{tcolorbox}
\section[\texorpdfstring{\te{కోర్టు} (kōrṭu)}{కోర్టు (kōrṭu)}]{\te{కోర్టు} (kōrṭu) — a court (of law)}

\label{lesson:TE-C256-kortu}



\begin{quote}
Before the new one: say the Telugu for the outskirts, then the Telugu for a warehouse.
\end{quote}

\subsection*{You'll want to know: \te{కోర్టు}}
\textbf{\te{కోర్టు}} — \emph{kōrṭu} — \textquotedblleft{}a court (of law)\textquotedblright{}.

\begin{grammarlens}[title={What you've built}]
One more word for this chapter.
\end{grammarlens}

\subsection*{Your turn}
\begin{itemize}
\raggedright
  \item \emph{Say it:} \emph{kōrṭu}
  \item \emph{Say it:} \emph{kōrṭu}, once more
  \item \emph{Say it:} say \emph{kōrṭu}
  \item \emph{Recall:} say the Telugu for the outskirts, then the Telugu for a warehouse, then say \emph{kōrṭu} again
\end{itemize}

\begin{tcolorbox}[breakable,colback=teal!4,colframe=teal!35!black,title={Before you move on}]
What does \te{కోర్టు} mean? (\textquotedblleft{}A court (of law)\textquotedblright{}.) Say it once more.
\end{tcolorbox}
